% 来源：OEBPS/Text/01_24.xhtml（OPF spine 顺序；仅第一册）
\cleardoublepage
\BookEntry[阿瑟·赫尔曼《苏格兰：现代世界文明的起点》]{阿瑟·赫尔曼《苏格兰：现代世界文明的起点》}

\BookCover{images/书单11苏格兰：现代世界文明的起点.jpg}{[美] 阿瑟·赫尔曼著| 启蒙编译所译| 上海社会科学院出版社2016年版}

\begin{BookSummaryBox}
\textbf{民主、科学与市场经济等现代世界的基石，何以诞生在18世纪的苏格兰？}
\par
\end{BookSummaryBox}

\subsubsection*{这本书谈什么？}

17世纪末，苏格兰还是欧洲最贫穷、最落后的国家之一，却奇迹般地在18世纪中叶迸发出惊人的能量，出现了经验主义哲学家休谟、自由经济的鼻祖亚当·斯密、影响欧洲浪漫主义深远的斯科特爵士、改良蒸汽机的瓦特等近代文明史中的巨人的身影。今日观之，举凡政治、法律、教育、科学、医学、建筑各领域，当时的苏格兰都是人才辈出。苏格兰不仅成为大不列颠帝国的中流砥柱，更随着“大西洋时代”的脚步，将其文明成就远播欧洲之外，启动全世界朝现代文明转向。
\par

地处偏远的蕞尔小国苏格兰，如何成为驱动现代世界运转的引擎？这本书将为读者解开这个连许多苏格兰人自己也无法回答的问题。
\par

\subsubsection*{芒格为什么谈起这本书？}

在西科金融公司2002财年的年会上，芒格说：我还推荐《苏格兰：现代世界文明的起点》这本书。这本书告诉了我们一小群杰出人物的成功故事。
\par

\subsubsection*{《苏格兰：现代世界文明的起点》作者阿瑟·赫尔曼如是说}

“道德判断（这是错的，那是对的）是人类的天赋，却不同于诸如判断距离远近的各种判断力。它会透过我们的情感加以呈现。爱是最重要的表达方式，尤其是对他人的爱，这是所有道德行为的出发点。”
\par

“这个文化与物质皆落后的国家，在思想前瞻的苏格兰人领导之下，即将展现另一番风貌。而当他们这么做的时候，世界也随之改观。到了18世纪末，象征现代世界的体制、理念、心态、习俗等特质，在苏格兰人手中已大致完备。苏格兰佬在全球披荆斩棘，为人类历史开展新的纪元；而所谓‘人类历史’的观念，也多半源自苏格兰的创见。”
\par

“苏格兰历史观的基础是强调进步。一如个人，社会本身也会与时俱进。社会需要新的技能、心态，对于个人能力以及个人应有何等自由，也要有新的理解。苏格兰人向世人揭示了衡量进步的指标之一——离过去的时代多远。以今衡古，而非以古论今。”
\par
